%========================================================
% CHAPTER 4 — Quantitative Characterization of Bipartite Quantum States
%========================================================

\chapter{Quantitative Characterization of Bipartite Quantum States}

%========================================================
% SECTION 4.1 — Reduced Density Operators as Local Diagnostics
%========================================================

\section{Reduced Density Operators as Local Diagnostics}
\label{sec:reduced_density_operators_diagnostics}

Reduced density operators were introduced in Section~\ref{subsec:reduced_density_operators} as the local states obtained from the global bipartite density operator by partial tracing over the complementary subsystem. In the present chapter, they are used as diagnostic tools for characterizing the local properties of each photon.

\medskip

For a two-photon state \( \rho_{AB} \), the reduced states are

\[
\rho_A
=
\mathrm{Tr}_B(\rho_{AB}),
\qquad
\rho_B
=
\mathrm{Tr}_A(\rho_{AB}).
\]

They determine all expectation values of observables acting locally on the corresponding subsystem. For instance, for an observable \( O_A \) acting only on arm \( A \),

\[
\langle O_A \rangle
=
\mathrm{Tr}(\rho_A O_A).
\]

Thus, \( \rho_A \) and \( \rho_B \) represent the complete local information accessible through measurements performed independently on the two optical arms.

\medskip

For Pauli observables, the local expectation values are

\begin{align}
\langle \sigma_x \rangle_A
&=
\mathrm{Tr}(\rho_A \sigma_x),
\nonumber
\\
\langle \sigma_y \rangle_A
&=
\mathrm{Tr}(\rho_A \sigma_y),
\nonumber
\\
\langle \sigma_z \rangle_A
&=
\mathrm{Tr}(\rho_A \sigma_z).
\label{eq:local_pauli_expectations_A}
\end{align}

Analogous expressions hold for subsystem \( B \). These quantities define the local Bloch-vector components of each reduced state and quantify the degree of local polarization.

\medskip

For the Bell state

\[
|\Psi^+\rangle
=
\frac{1}{\sqrt{2}}
\left(
|01\rangle + |10\rangle
\right),
\]

one obtains

\[
\rho_A
=
\rho_B
=
\frac{1}{2}I.
\]

Consequently,

\[
\langle \sigma_x \rangle_A
=
\langle \sigma_y \rangle_A
=
\langle \sigma_z \rangle_A
=
0,
\]

and similarly for subsystem \( B \). Each photon is therefore locally unpolarized, even though the global two-photon state is pure and maximally entangled.

\medskip

The same local behavior holds for the reduced phase model introduced in Section~\ref{sec:reduced_phase_model}. The state

\[
|\psi(\theta)\rangle
=
\frac{1}{\sqrt{2}}
\left(
|01\rangle + e^{i\theta}|10\rangle
\right)
\]

has reduced density operators

\[
\rho_A(\theta)
=
\rho_B(\theta)
=
\frac{1}{2}I.
\]

The relative phase \( \theta \) changes the global coherence structure and the two-photon correlations, but it does not modify the local reduced states. Therefore, local polarization measurements alone cannot detect the deterministic phase shift.

\medskip

By linearity of the partial trace, the same conclusion extends to the random-phase ensemble introduced in Section~\ref{sec:pure_and_density_operator_descriptions}. If

\[
\rho_{AB}^{(\mathrm{ens})}
=
\sum_k q_k
|\psi(\theta_k)\rangle
\langle \psi(\theta_k)|,
\qquad
\sum_k q_k = 1,
\]

then

\[
\rho_A^{(\mathrm{ens})}
=
\rho_B^{(\mathrm{ens})}
=
\frac{1}{2}I.
\]

Thus, in the reduced phase model and in its statistical extension, the local states remain maximally mixed. The effect of the phase distribution appears at the global level, through the decay of coherences and correlations, rather than through a change in local polarization.

\medskip

This observation is important for the interpretation of entanglement. The same local density operator

\[
\rho_A
=
\frac{1}{2}I
\]

can arise from physically different situations. It may describe a classical statistical mixture,

\[
\rho^{(\mathrm{class})}
=
\frac{1}{2}
|0\rangle\langle 0|
+
\frac{1}{2}
|1\rangle\langle 1|,
\]

or it may arise as the reduced state of an entangled bipartite system,

\[
\rho_A^{(\mathrm{ent})}
=
\mathrm{Tr}_B
\left(
|\Psi^+\rangle\langle\Psi^+|
\right)
=
\frac{1}{2}I.
\]

These two cases are indistinguishable by local measurements alone. Their difference is visible only at the global level, where one must analyze global purity, bipartite correlations, and entanglement measures.

\medskip

For this reason, reduced density operators are not sufficient by themselves to characterize the full bipartite state. They provide local information, while quantities such as global purity, fidelity, concurrence, tangle, entanglement of formation, and CHSH parameters are required to quantify the global coherence, entanglement, and nonlocality of the two-photon system.

%========================================================
% SECTION 4.2 — Purity and Mixedness
%========================================================

\section{Purity and Mixedness}

Having introduced reduced density operators in Section~4.1, we now require a quantitative criterion to distinguish whether a quantum state is pure or mixed.

\medskip

A state described by a density operator \( \rho \) is pure if it can be written as

\begin{equation}
\rho
=
|\psi\rangle\langle\psi|,
\label{eq:pure_density_operator}
\end{equation}

for some normalized vector \( |\psi\rangle \). Otherwise, it is mixed, meaning that it admits an ensemble representation

\begin{equation}
\rho
=
\sum_k
p_k
|\psi_k\rangle\langle\psi_k|,
\qquad
p_k \ge 0,
\qquad
\sum_k p_k = 1.
\label{eq:mixed_density_operator}
\end{equation}

\medskip

A convenient quantitative measure of mixedness is the purity, defined as

\begin{equation}
\gamma(\rho)
=
\mathrm{Tr}(\rho^2).
\label{eq:purity_definition}
\end{equation}

This quantity satisfies

\begin{equation}
\gamma(\rho)=1
\Longleftrightarrow
\rho \text{ is pure},
\qquad
\gamma(\rho)<1
\Longleftrightarrow
\rho \text{ is mixed}.
\label{eq:purity_pure_mixed_condition}
\end{equation}

\medskip

If \( \rho = \sum_i \lambda_i |i\rangle\langle i| \) is its spectral decomposition, then

\begin{equation}
\gamma(\rho)
=
\sum_i \lambda_i^2,
\label{eq:purity_spectral_decomposition}
\end{equation}

which shows that purity measures the concentration of the eigenvalue distribution.

More generally, for a Hilbert space of dimension \( d \),

\begin{equation}
\frac{1}{d}
\le
\gamma(\rho)
\le
1,
\label{eq:purity_bounds}
\end{equation}

with the lower bound attained by the maximally mixed state \( \rho = I/d \).

\medskip

For single-qubit systems, purity admits a simple geometric interpretation. Any density operator can be written as

\begin{equation}
\rho
=
\frac{1}{2}
\left(
I + \vec{r}\cdot\vec{\sigma}
\right),
\label{eq:single_qubit_bloch_representation}
\end{equation}

where \( \vec{r} \in \mathbb{R}^3 \) is the Bloch vector. One then finds

\begin{equation}
\gamma(\rho)
=
\frac{1+\|\vec{r}\|^2}{2}.
\label{eq:single_qubit_purity_bloch}
\end{equation}

Thus,

\[
\|\vec{r}\| = 1
\Rightarrow
\rho \text{ is pure},
\qquad
\|\vec{r}\| = 0
\Rightarrow
\rho = \frac{I}{2}.
\]

\medskip

Mixedness does not uniquely determine the physical origin of a state. As discussed in Section~4.1, it may arise either from classical statistical uncertainty or from entanglement with an unobserved subsystem.

\medskip

A paradigmatic example is provided by the Bell state

\[
|\Psi^+\rangle
=
\frac{|01\rangle + |10\rangle}{\sqrt{2}},
\]

whose global density operator

\[
\rho_{AB}
=
|\Psi^+\rangle\langle\Psi^+|
\]

is pure,

\[
\gamma(\rho_{AB}) = 1,
\]

while the reduced states satisfy

\[
\rho_A
=
\rho_B
=
\frac{I}{2},
\qquad
\gamma(\rho_A)
=
\gamma(\rho_B)
=
\frac{1}{2}.
\]

This illustrates that a globally pure entangled state may generate locally maximally mixed subsystems.

\medskip

In the present model, the phase-evolved state introduced in Chapter~3,

\[
|\psi(\theta)\rangle
=
\frac{|01\rangle + e^{i\theta}|10\rangle}{\sqrt{2}},
\]

remains pure for all values of \( \theta \),

\begin{equation}
\gamma(\rho(\theta)) = 1,
\label{eq:phase_state_global_purity}
\end{equation}

while the reduced states remain maximally mixed,

\[
\rho_A(\theta)
=
\rho_B(\theta)
=
\frac{I}{2}.
\]

Thus, relative phase shifts modify correlation observables without affecting either global purity or local mixedness.

\medskip

When extending the model to include stochastic or environmental effects, the situation changes. For the random-phase ensemble, the global state depends on the coherence parameter

\begin{equation}
\mu
=
\langle e^{i\theta} \rangle,
\label{eq:coherence_parameter_mu}
\end{equation}

and its purity is

\begin{equation}
\gamma
\left(
\rho_{AB}^{(\mathrm{ens})}
\right)
=
\frac{1+|\mu|^2}{2}.
\label{eq:random_phase_ensemble_purity_mu}
\end{equation}

Equivalently,

\begin{equation}
\gamma
\left(
\rho_{AB}^{(\mathrm{ens})}
\right)
=
\frac{
1
+
\langle \cos\theta \rangle^2
+
\langle \sin\theta \rangle^2
}{2}.
\label{eq:random_phase_ensemble_purity_trigonometric}
\end{equation}

This shows that purity directly quantifies the loss of phase coherence in the ensemble: the state is pure for \( |\mu|=1 \) and becomes mixed as \( |\mu| \) decreases.

\medskip

Purity therefore provides a fundamental diagnostic tool for assessing coherence loss in the simulator, complementing correlation observables and preparing the ground for fidelity and entanglement measures.

%========================================================
% SECTION 4.3 — Fidelity with Respect to Bell States
%========================================================

\section{Fidelity with Respect to Bell States}

Having introduced reduced density operators and purity as measures of local information and mixedness, we now require a quantitative tool to evaluate how close a given quantum state is to a reference target state. This role is fulfilled by the notion of fidelity.

\medskip

In general, the fidelity between two quantum states described by density operators \( \rho \) and \( \sigma \) is defined as

\begin{equation}
F(\rho,\sigma)
=
\left(
\mathrm{Tr}
\sqrt{
\sqrt{\rho}\,
\sigma\,
\sqrt{\rho}
}
\right)^2,
\label{eq:fidelity_general_definition}
\end{equation}

and satisfies

\[
0
\le
F(\rho,\sigma)
\le
1,
\]

with equality to \(1\) if and only if \( \rho = \sigma \).

\medskip

In the present context, the reference state is pure, \( \sigma = |\phi\rangle\langle\phi| \), and the expression simplifies to

\begin{equation}
F(\rho,|\phi\rangle)
=
\langle \phi | \rho | \phi \rangle,
\label{eq:fidelity_pure_reference}
\end{equation}

which admits a direct interpretation as the probability of projecting onto \( |\phi\rangle \).

\medskip

We focus on the Bell state introduced in Chapter~3,

\[
|\Psi^+\rangle
=
\frac{|01\rangle + |10\rangle}{\sqrt{2}},
\]

and define

\begin{equation}
F_{\Psi^+}(\rho)
=
\langle \Psi^+ | \rho | \Psi^+ \rangle.
\label{eq:fidelity_psi_plus_definition}
\end{equation}

\medskip

For the phase-evolved state,

\[
|\psi(\theta)\rangle
=
\frac{|01\rangle + e^{i\theta}|10\rangle}{\sqrt{2}},
\qquad
\rho(\theta)
=
|\psi(\theta)\rangle\langle\psi(\theta)|,
\]

one obtains

\[
F_{\Psi^+}(\rho(\theta))
=
|\langle \Psi^+|\psi(\theta)\rangle|^2,
\]

with

\[
\langle \Psi^+|\psi(\theta)\rangle
=
\frac{1}{2}
\left(
1+e^{i\theta}
\right),
\]

leading to

\begin{equation}
F_{\Psi^+}(\rho(\theta))
=
\frac{1+\cos\theta}{2}.
\label{eq:phase_state_fidelity_psi_plus}
\end{equation}

\medskip

Thus, the relative phase modifies the overlap with the reference Bell state, even though the global state remains pure. In particular,

\[
F_{\Psi^+}=1
\quad
\text{for }
\theta=0,
\qquad
F_{\Psi^+}=0
\quad
\text{for }
\theta=\pi.
\]

\medskip

For ensemble states,

\[
\rho_{AB}^{(\mathrm{ens})}
=
\frac{1}{N}
\sum_{k=1}^{N}
|\psi(\theta_k)\rangle\langle\psi(\theta_k)|,
\]

linearity yields

\[
F_{\Psi^+}
\left(
\rho_{AB}^{(\mathrm{ens})}
\right)
=
\frac{1}{N}
\sum_{k=1}^{N}
F_{\Psi^+}
\left(
\rho(\theta_k)
\right).
\]

Introducing the coherence parameter \( \mu = \langle e^{i\theta}\rangle \), one finds

\begin{equation}
F_{\Psi^+}
=
\frac{1}{2}
\left(
1+\mathrm{Re}(\mu)
\right)
=
\frac{1}{2}
\left(
1+\langle \cos\theta\rangle
\right).
\label{eq:ensemble_fidelity_psi_plus}
\end{equation}

\medskip

In particular, for a uniform distribution over \( [0,2\pi] \),

\[
F_{\Psi^+}
=
\frac{1}{2}.
\]

\medskip

Fidelity therefore provides a quantitative measure of how closely a given state approximates the target Bell state. Unlike purity, which characterizes mixedness, fidelity is sensitive to coherent deformations of the state and complements correlation observables in the analysis of the system.

%========================================================
% SECTION 4.4 — Concurrence as an Entanglement Measure
%========================================================

\section{Concurrence as an Entanglement Measure}

Reduced density operators, purity, and fidelity provide complementary information about a bipartite quantum state, but none of them directly quantifies entanglement. In particular, fidelity with respect to a Bell state may vary under local unitary evolution even when the amount of entanglement remains unchanged. A genuine entanglement measure is therefore required.

\medskip

For two-qubit systems, a standard measure is the concurrence, defined such that

\[
0 \le C \le 1,
\]

where \( C=0 \) corresponds to separable states and \( C=1 \) to maximally entangled states.

\medskip

For a pure state

\[
|\psi\rangle
=
a|00\rangle
+
b|01\rangle
+
c|10\rangle
+
d|11\rangle,
\]

the concurrence is

\begin{equation}
C(|\psi\rangle)
=
2|ad-bc|.
\label{eq:pure_state_concurrence}
\end{equation}

This expression vanishes for separable states and reaches unity for Bell states. For example,

\[
|\Psi^+\rangle
=
\frac{|01\rangle + |10\rangle}{\sqrt{2}}
\quad
\Rightarrow
\quad
C=1.
\]

\medskip

For the phase-evolved state introduced in Chapter~3,

\[
|\psi(\theta)\rangle
=
\frac{|01\rangle + e^{i\theta}|10\rangle}{\sqrt{2}},
\]

one finds

\begin{equation}
C(|\psi(\theta)\rangle)
=
1,
\label{eq:phase_state_concurrence}
\end{equation}

showing that the relative phase modifies fidelity but leaves entanglement unchanged. This reflects the invariance of bipartite entanglement under local unitary transformations.

\medskip

For pure bipartite states, concurrence is directly related to the reduced density operators:

\begin{equation}
C(|\psi\rangle)
=
2\sqrt{\det(\rho_A)}
=
\sqrt{
2\left(
1-\gamma(\rho_A)
\right)
},
\qquad
\rho_A
=
\mathrm{Tr}_B
\left(
|\psi\rangle\langle\psi|
\right).
\label{eq:concurrence_reduced_density_relation}
\end{equation}

Thus, entanglement is encoded in the mixedness of the reduced states: separable states yield pure reductions, while maximally entangled states yield maximally mixed reductions.

\medskip
s
For mixed states \( \rho \in \mathbb{C}^{4 \times 4} \), concurrence is defined through the Wootters spin-flip construction \cite{wootters_entanglement_1998}. The associated spin-flipped density operator is

\begin{equation}
\tilde{\rho}
=
(\sigma_y \otimes \sigma_y)
\rho^*
(\sigma_y \otimes \sigma_y).
\label{eq:spin_flipped_density_matrix}
\end{equation}

Let \( \lambda_1 \ge \lambda_2 \ge \lambda_3 \ge \lambda_4 \ge 0 \) be the square roots of the eigenvalues of \( \rho\tilde{\rho} \). Then

\begin{equation}
C(\rho)
=
\max
\left(
0,
\lambda_1-\lambda_2-\lambda_3-\lambda_4
\right).
\label{eq:mixed_state_concurrence}
\end{equation}

This expression applies to arbitrary two-qubit states and reduces to the pure-state formula in the appropriate limit.

\medskip

In the present model, mixed states arise from ensemble averaging over random phases,

\[
\rho_{AB}^{(\mathrm{ens})}
=
\frac{1}{N}
\sum_{k=1}^{N}
|\psi(\theta_k)\rangle\langle\psi(\theta_k)|,
\]

which depend on the coherence parameter

\[
\mu
=
\langle e^{i\theta}\rangle.
\]

In the computational basis, the ensemble state reads

\begin{equation}
\rho_{AB}^{(\mathrm{ens})}
=
\frac{1}{2}|01\rangle\langle 01|
+
\frac{1}{2}|10\rangle\langle 10|
+
\frac{\mu}{2}|10\rangle\langle 01|
+
\frac{\mu^*}{2}|01\rangle\langle 10|.
\label{eq:random_phase_ensemble_density_explicit}
\end{equation}

For this family of states, the concurrence is

\begin{equation}
C
\left(
\rho_{AB}^{(\mathrm{ens})}
\right)
=
|\mu|.
\label{eq:random_phase_ensemble_concurrence}
\end{equation}

\medskip

This result shows that the same coherence parameter controlling purity and fidelity also governs entanglement degradation:

\[
\gamma
\left(
\rho_{AB}^{(\mathrm{ens})}
\right)
=
\frac{1+|\mu|^2}{2},
\qquad
F_{\Psi^+}
=
\frac{1+\mathrm{Re}(\mu)}{2}.
\]

Thus, the reduced phase model provides a unified description in which coherence, mixedness, and entanglement are all controlled by a single parameter.

%========================================================
% SECTION 4.5 — Tangle
%========================================================

\section{Tangle}

While concurrence provides a direct quantitative measure of bipartite entanglement, it is often useful to introduce an associated quadratic quantity known as the \emph{tangle}.

For a two-qubit state \( \rho \), the tangle is defined as

\begin{equation}
\tau(\rho)
=
C(\rho)^2,
\label{eq:tangle_definition}
\end{equation}

where \( C(\rho) \) denotes the concurrence.

\medskip

Since concurrence satisfies

\[
0
\leq
C(\rho)
\leq
1,
\]

the tangle also satisfies

\[
0
\leq
\tau(\rho)
\leq
1.
\]

The tangle therefore vanishes for separable states and reaches unity for maximally entangled Bell states.

\medskip

For pure bipartite states, the tangle admits an equivalent representation in terms of the reduced density operator. If

\[
|\psi\rangle
\in
\mathcal{H}_A
\otimes
\mathcal{H}_B,
\]

then

\begin{equation}
\tau
=
4
\det(\rho_A),
\label{eq:tangle_reduced_density_operator}
\end{equation}

where

\[
\rho_A
=
\mathrm{Tr}_B
\left(
|\psi\rangle
\langle\psi|
\right).
\]

This expression shows directly that the tangle quantifies the degree of non-purity of the reduced subsystem induced by bipartite entanglement.

\medskip

For the Bell state \( |\Psi^+\rangle \), one obtains

\[
C
=
1,
\qquad
\tau
=
1,
\]

while for separable product states,

\[
C
=
0,
\qquad
\tau
=
0.
\]

\medskip

Compared to concurrence, the quadratic dependence of the tangle emphasizes the suppression of weak entanglement contributions. Consequently, the tangle often exhibits a sharper degradation under noisy evolution.

\medskip

The tangle also plays an important role in the study of multipartite entanglement sharing and monogamy relations, where squared concurrence naturally appears in entanglement-balance inequalities.

\medskip

Within the present work, the tangle is used as an auxiliary entanglement metric complementing concurrence and providing an alternative characterization of entanglement degradation under effective fiber propagation.

%========================================================
% SECTION 4.6 — Entanglement of Formation
%========================================================

\section{Entanglement of Formation}

The concurrence and tangle quantify the amount of bipartite entanglement present in a quantum state. However, they do not directly describe the operational cost required to create that entanglement.

For this purpose, a central quantity in quantum information theory is the \emph{entanglement of formation}.

\medskip

The entanglement of formation measures the minimal amount of bipartite entanglement required, on average, to prepare a given mixed state using pure entangled states and local operations with classical communication.

\medskip

For arbitrary quantum systems, the computation of entanglement of formation is generally difficult. Remarkably, however, Wootters showed that the entanglement of formation of arbitrary two-qubit states can be expressed analytically in terms of concurrence \cite{wootters_entanglement_1998}.

\medskip

For a two-qubit density operator \( \rho \), the entanglement of formation is defined as

\begin{equation}
E_F(\rho)
=
h
\left(
\frac{
1
+
\sqrt{1-C(\rho)^2}
}{2}
\right),
\label{eq:eof_definition}
\end{equation}

where \( C(\rho) \) denotes the concurrence and

\begin{equation}
h(x)
=
-
x \log_2 x
-
(1-x)\log_2(1-x)
\label{eq:binary_entropy_function}
\end{equation}

is the binary entropy function.

\medskip

Since

\[
0
\leq
C(\rho)
\leq
1,
\]

the entanglement of formation satisfies

\[
0
\leq
E_F(\rho)
\leq
1.
\]

In particular:

\[
E_F
=
0
\]

for separable states, while

\[
E_F
=
1
\]

for maximally entangled Bell states.

\medskip

Unlike concurrence, the entanglement of formation depends nonlinearly on the entanglement content of the state. As a consequence, small reductions in concurrence may produce strongly nonlinear reductions in \( E_F \).

\medskip

This nonlinear behavior is physically significant because it reflects the effective loss of entanglement resources available for quantum-information tasks such as teleportation, entanglement distribution, and quantum communication protocols.

\medskip

For the Bell state \( |\Psi^+\rangle \), one obtains

\[
C
=
1,
\qquad
E_F
=
1,
\]

while for separable states,

\[
C
=
0,
\qquad
E_F
=
0.
\]

\medskip

Within the present work, the entanglement of formation provides an operational complement to concurrence and tangle, allowing the degradation of entanglement under effective fiber propagation to be interpreted not only geometrically, but also in terms of entanglement-resource loss.

%========================================================
% SECTION 4.7 — Bell Inequalities and CHSH Nonlocality
%========================================================

\section{Bell Inequalities and CHSH Nonlocality}

The correlation structure introduced in the previous sections provides a complete operational description of measurement outcomes for arbitrary local observables. This framework makes it possible to formulate quantitative criteria that distinguish correlations compatible with classical physics from those that are intrinsically quantum.

\medskip

The central question addressed by Bell inequalities is the following: \emph{can the observed correlations between spatially separated subsystems be explained by a classical model based on shared hidden information?}

\medskip

In a classical description based on local hidden variables, it is assumed that the outcomes of measurements performed on subsystems \( A \) and \( B \) are determined by a shared underlying variable \( \lambda \), drawn from a probability distribution \( p(\lambda) \). This variable represents all pre-existing information that could influence the measurement results.

\medskip

For each realization of \( \lambda \), the outcome of a measurement on subsystem \( A \) along direction \( \mathbf{a} \) is described by a deterministic response function

\[
A(\mathbf{a},\lambda) \in \{-1,+1\},
\]

and similarly for subsystem \( B \),

\[
B(\mathbf{b},\lambda) \in \{-1,+1\}.
\]

\medskip

The key physical assumption is locality: the outcome at each subsystem depends only on the local measurement setting and on the shared variable \( \lambda \), but not on the measurement choice at the distant subsystem.

Since \( \lambda \) is not directly accessible, observable quantities correspond to averages over its distribution. The classical correlation function is therefore \cite{bell_epr_paradox_1964}

\begin{equation}
E(\mathbf{a},\mathbf{b})
=
\int d\lambda\,
p(\lambda)\,
A(\mathbf{a},\lambda)\,
B(\mathbf{b},\lambda),
\qquad
p(\lambda) \ge 0,
\qquad
\int d\lambda\,p(\lambda)=1.
\label{eq:local_hidden_variable_correlation}
\end{equation}

This expression represents the most general form of bipartite correlations compatible with a local hidden-variable model.

\medskip

The most widely used form of a Bell inequality is the Clauser--Horne--Shimony--Holt (CHSH) inequality. It involves four measurement settings:

\[
\mathbf{a}_1,\mathbf{a}_2
\quad
\text{for subsystem } A,
\qquad
\mathbf{b}_1,\mathbf{b}_2
\quad
\text{for subsystem } B.
\]

The CHSH parameter is defined as

\begin{equation}
S
=
E(\mathbf{a}_1,\mathbf{b}_1)
+
E(\mathbf{a}_1,\mathbf{b}_2)
+
E(\mathbf{a}_2,\mathbf{b}_1)
-
E(\mathbf{a}_2,\mathbf{b}_2).
\label{eq:chsh_parameter}
\end{equation}

\medskip

\paragraph{Classical bound.}

Under the assumptions of locality and realism, the CHSH parameter satisfies

\begin{equation}
|S| \le 2.
\label{eq:chsh_classical_bound}
\end{equation}

This bound reflects a fundamental limitation of classical correlations: although each term in \( S \) can individually reach its maximal value, their combination is constrained by the requirement that all outcomes originate from the same underlying variable \( \lambda \).

\medskip

\paragraph{Quantum prediction.}

In quantum mechanics, correlations arise from the structure of the joint state rather than from pre-existing local variables. The correlation function is given by Eq.~\eqref{eq:axis_correlation_trace}.

Due to the non-commuting nature of quantum observables and the presence of entanglement, the CHSH parameter can exceed the classical bound. The maximal value allowed by quantum mechanics is

\begin{equation}
|S| \le 2\sqrt{2},
\label{eq:tsirelson_bound}
\end{equation}

known as the Tsirelson bound.

This violation does not imply faster-than-light communication. Rather, it implies that the correlations cannot be reproduced by any model based on local hidden variables.

\medskip

\paragraph{Tensor formulation.}

Using the correlation tensor representation, the CHSH parameter can be expressed as

\begin{equation}
S
=
\mathbf{a}_1^{\mathsf T} T \mathbf{b}_1
+
\mathbf{a}_1^{\mathsf T} T \mathbf{b}_2
+
\mathbf{a}_2^{\mathsf T} T \mathbf{b}_1
-
\mathbf{a}_2^{\mathsf T} T \mathbf{b}_2.
\label{eq:chsh_tensor_form}
\end{equation}

This formulation makes explicit that the CHSH inequality depends only on the correlation tensor and on the choice of measurement directions. In particular, it shows that the observed value of \( S \) is sensitive not only to the state, but also to the alignment between the measurement axes and the principal directions of the tensor.

\medskip

A particularly important result is that the maximal violation achievable for a given state can be computed directly from the tensor. Defining

\begin{equation}
U = T^{\mathsf T}T,
\label{eq:horodecki_matrix}
\end{equation}

and denoting by \( \lambda_1 \) and \( \lambda_2 \) the two largest eigenvalues of \( U \), the maximal CHSH value is

\begin{equation}
S_{\max}
=
2\sqrt{\lambda_1+\lambda_2}.
\label{eq:chsh_smax_horodecki}
\end{equation}

\medskip

\paragraph{Interpretation.}

A violation of the CHSH inequality, i.e. \( |S| > 2 \), indicates the presence of correlations that cannot be explained by any local hidden-variable model. Such correlations are referred to as nonlocal.

\medskip

Nonlocality should be understood as a property of correlations rather than of individual subsystems: each subsystem alone may appear completely random, while the joint statistics reveal strong structured correlations.

\medskip

It is important to distinguish nonlocality from entanglement. While all maximally entangled pure states achieve maximal quantum violation for suitable measurement axes, not all entangled states violate the CHSH inequality. In particular, sufficiently mixed states may remain within the classical bound despite being entangled.

\medskip

\paragraph{Operational implications.}

The CHSH parameter provides a direct bridge between the theoretical description of correlations and experimentally accessible quantities. In practice, it is evaluated by performing measurements along four selected pairs of local axes and combining the resulting correlation values.

\medskip

The correlation tensor entering this construction can be experimentally reconstructed from a set of bipartite expectation values in the Pauli basis, in a procedure closely related to quantum state tomography. In this sense, the evaluation of CHSH inequalities can be viewed as a specific post-processing of tomographically accessible data.

\medskip

Two complementary perspectives are particularly relevant:
\begin{itemize}
    \item evaluation of \( S \) for fixed measurement settings, corresponding to a specific experimental configuration,
    \item computation of \( S_{\max} \), which quantifies the maximal nonlocality extractable from the state.
\end{itemize}

\medskip

This separation will play a central role in the analysis of fiber-induced effects, where transformations may alter the observed correlations without necessarily reducing the underlying nonlocality.

\medskip

%========================================================
% SECTION 4.8 — Coherence, Entanglement, and Nonlocality in the Random-Phase Model
%========================================================

\section{Coherence, Entanglement, and Nonlocality in the Random-Phase Model}

The previous sections introduced several complementary quantities used to characterize bipartite quantum states:

\begin{itemize}
    \item purity, quantifying mixedness;
    \item fidelity, quantifying overlap with a reference state;
    \item concurrence and tangle, quantifying bipartite entanglement;
    \item entanglement of formation, quantifying entanglement resources;
    \item CHSH nonlocality, quantifying violation of local hidden-variable models.
\end{itemize}

\medskip

Although these quantities are related, they characterize fundamentally different physical properties of the state.

\medskip

Purity measures how close the density operator is to a pure state, independently of whether the state is entangled or separable.

Fidelity measures proximity to a chosen target state and therefore depends explicitly on the selected reference.

Concurrence and tangle quantify bipartite entanglement, while the entanglement of formation provides an operational interpretation of the entanglement content in terms of preparation resources.

Finally, CHSH inequalities probe nonlocal correlations that cannot be reproduced by any local hidden-variable model.

\medskip

These distinctions are particularly important because entanglement and nonlocality are not equivalent concepts.

\medskip

All pure maximally entangled Bell states violate the CHSH inequality maximally for suitable measurement settings. However, in general:

\[
\text{entanglement}
\neq
\text{nonlocality}.
\]

In particular, sufficiently mixed entangled states may fail to violate any CHSH inequality despite possessing nonzero concurrence and nonzero entanglement of formation.

\medskip

The random-phase ensemble considered in this work provides a particularly useful example because the state properties are governed by a single coherence parameter

\begin{equation}
\mu
=
\langle e^{i\theta}\rangle
=
\langle \cos\theta\rangle
+
i\langle \sin\theta\rangle.
\label{eq:coherence_parameter_trigonometric_updated}
\end{equation}

For this family of states, the relevant quantities become

\begin{equation}
\gamma(\rho)
=
\frac{1+|\mu|^2}{2},
\label{eq:purity_mu_relation}
\end{equation}

\begin{equation}
C(\rho)
=
|\mu|,
\qquad
\tau(\rho)
=
|\mu|^2,
\label{eq:concurrence_tangle_mu_relation}
\end{equation}

\begin{equation}
F_{\Psi^+}
=
\frac{1+\mathrm{Re}(\mu)}{2},
\label{eq:fidelity_mu_relation}
\end{equation}

and

\begin{equation}
E_F(\rho)
=
h
\left(
\frac{
1+\sqrt{1-|\mu|^2}
}{2}
\right),
\label{eq:eof_mu_relation}
\end{equation}

where \( h(x) \) denotes the binary entropy function introduced previously.

\medskip

These expressions reveal the different sensitivities of the various quantities.

\begin{itemize}
    \item Purity depends quadratically on \( |\mu| \).
    \item Concurrence depends linearly on \( |\mu| \).
    \item Tangle depends quadratically on \( |\mu| \).
    \item Entanglement of formation depends nonlinearly on \( |\mu| \).
    \item Fidelity depends only on the real part \( \mathrm{Re}(\mu) \).
\end{itemize}

\medskip

Consequently, two states may possess the same concurrence while having different fidelities, or may possess nonzero entanglement while exhibiting no observable CHSH violation.

\medskip

This distinction becomes particularly evident in the deterministic phase model. In that case,

\[
|\mu|
=
1
\]

for all phase values \( \theta \), because the evolution remains purely unitary. Therefore,

\[
C
=
1,
\qquad
\tau
=
1,
\qquad
E_F
=
1,
\]

while the fidelity varies continuously with the phase:

\[
F_{\Psi^+}
=
\frac{1+\cos\theta}{2}.
\]

Thus, changes in fidelity do not necessarily correspond to degradation of entanglement.

\medskip

By contrast, ensemble averaging reduces the magnitude \( |\mu| \), producing genuine decoherence and entanglement degradation. In this regime:

\begin{itemize}
    \item the global purity decreases;
    \item concurrence decreases;
    \item tangle decreases more rapidly due to its quadratic dependence;
    \item entanglement of formation decreases nonlinearly;
    \item CHSH nonlocality may eventually disappear.
\end{itemize}

\medskip

A useful geometric interpretation is obtained by representing \( \mu \) in the complex plane.

\begin{itemize}
    \item The concurrence is proportional to the distance of \( \mu \) from the origin.
    \item The tangle depends on the squared distance.
    \item The fidelity corresponds to the projection onto the real axis.
\end{itemize}

Deterministic phase evolution moves \( \mu \) along the unit circle without changing its magnitude, while ensemble averaging contracts \( \mu \) toward the origin.

\medskip

The limiting regimes illustrate these behaviors clearly:

\begin{equation}
|\mu|=1
\Rightarrow
\gamma(\rho)=1,
\quad
C=1,
\quad
\tau=1,
\quad
E_F=1,
\label{eq:coherence_parameter_limit_pure}
\end{equation}

while

\begin{equation}
|\mu|=0
\Rightarrow
\gamma(\rho)=\frac{1}{2},
\quad
C=0,
\quad
\tau=0,
\quad
E_F=0.
\label{eq:coherence_parameter_limit_mixed}
\end{equation}

Intermediate values \( 0<|\mu|<1 \) correspond to partially coherent mixed states with nonzero but degraded entanglement.

\medskip

These relations are specific to the random-phase ensemble model considered in this work. For general two-qubit states, purity, fidelity, concurrence, entanglement of formation, and CHSH nonlocality are independent quantities and cannot generally be reduced to a single parameter.

\medskip

In summary, the coherence parameter \( \mu \) provides a unified description of the states analyzed throughout the present work. Its magnitude controls mixedness and entanglement, while its phase determines the orientation of observable correlations and the overlap with the reference Bell state.

The comparison between these quantities highlights an important conceptual hierarchy:

\[
\text{mixedness}
\quad
\longrightarrow
\quad
\text{entanglement}
\quad
\longrightarrow
\quad
\text{nonlocality}.
\]

This hierarchy will play a central role in the interpretation of the numerical experiments and in the analysis of effective fiber-induced degradation mechanisms.

%========================================================
% SECTION 4.9 — Computational Evaluation of State Metrics
%========================================================

\section{Computational Evaluation of State Metrics}

The quantities introduced throughout this chapter constitute the main diagnostic layer of the simulator. Their role is to translate the density-operator description of bipartite quantum states into quantitative indicators of mixedness, coherence, entanglement, and nonlocality.

\medskip

At the computational level, all metrics are evaluated directly from the numerically reconstructed density operator \( \rho_{AB} \). This provides a unified characterization framework applicable to deterministic phase evolution, statistical phase ensembles, depolarizing channels, and composite effective fiber models.

\medskip

Reduced density operators are obtained through partial-trace operations acting on the global state:

\[
\rho_A
=
\mathrm{Tr}_B(\rho_{AB}),
\qquad
\rho_B
=
\mathrm{Tr}_A(\rho_{AB}).
\]

For two-qubit systems represented in the computational basis, these operations return \(2\times2\) matrices that encode the effective local states of the two subsystems.

\medskip

The reduced states are then used to evaluate local purities, local expectation values, and entanglement-related quantities.

\medskip

\paragraph{Purity evaluation.}

Purity is computed numerically using the trace formula

\[
\gamma(\rho)
=
\mathrm{Tr}(\rho^2),
\]

which is applied both to the global density operator and to the reduced density operators.

\medskip

This quantity provides a direct computational diagnostic of coherence loss and mixedness generation. In particular:

\begin{itemize}
    \item pure states satisfy
    \[
    \gamma(\rho)=1,
    \]
    \item mixed states satisfy
    \[
    \gamma(\rho)<1,
    \]
    \item maximally mixed single-qubit states satisfy
    \[
    \gamma=\frac{1}{2}.
    \]
\end{itemize}

\medskip

\paragraph{Fidelity evaluation.}

Fidelity with respect to the reference Bell state \( |\Psi^+\rangle \) is computed using the pure-reference formula

\[
F_{\Psi^+}(\rho)
=
\langle\Psi^+|
\rho
|\Psi^+\rangle.
\]

Since the reference state is pure, this expression avoids the matrix square-root operations required by the general Uhlmann fidelity and therefore constitutes the natural implementation for the present simulator.

\medskip

\paragraph{Concurrence evaluation.}

Concurrence is evaluated using the standard two-qubit spin-flip construction. Defining

\[
\tilde{\rho}
=
(\sigma_y\otimes\sigma_y)
\rho^*
(\sigma_y\otimes\sigma_y),
\]

the simulator computes the eigenvalues of

\[
\rho\tilde{\rho},
\]

and evaluates the concurrence according to the mixed-state formula introduced previously.

\medskip

For pure states, the analytical expression

\[
C(|\psi\rangle)
=
2|ad-bc|
\]

provides an additional validation benchmark for the numerical implementation.

\medskip

\paragraph{Tangle and entanglement of formation.}

The tangle is evaluated directly from the concurrence through

\[
\tau(\rho)
=
C(\rho)^2,
\]

while the entanglement of formation is computed from

\[
E_F(\rho)
=
h
\left(
\frac{
1+\sqrt{1-C(\rho)^2}
}{2}
\right),
\]

where \( h(x) \) denotes the binary entropy function.

\medskip

Implementing both quantities through concurrence guarantees numerical consistency across the entanglement-diagnostic layer while avoiding redundant spectral decompositions.

\medskip

\paragraph{CHSH nonlocality evaluation.}

The CHSH parameter is evaluated using the correlation tensor representation introduced previously.

For fixed measurement settings, the simulator computes

\[
S
=
\mathbf{a}_1^{\mathsf T} T \mathbf{b}_1
+
\mathbf{a}_1^{\mathsf T} T \mathbf{b}_2
+
\mathbf{a}_2^{\mathsf T} T \mathbf{b}_1
-
\mathbf{a}_2^{\mathsf T} T \mathbf{b}_2.
\]

In addition, the maximal CHSH violation is evaluated through the Horodecki construction:

\[
S_{\max}
=
2\sqrt{\lambda_1+\lambda_2},
\]

where \( \lambda_1 \) and \( \lambda_2 \) denote the two largest eigenvalues of

\[
U=T^{\mathsf T}T.
\]

\medskip

This distinction allows the simulator to separate experimentally observed nonlocality for specific measurement settings from the maximal nonlocality intrinsically supported by the state.

\medskip

\paragraph{Unified diagnostic framework.}

All numerical experiments presented in this work employ the same set of computational metric routines. This ensures that deterministic evolution, statistical averaging, and effective noisy channels are characterized within a common and internally consistent framework.

\medskip

The computational objective is therefore not merely the evaluation of isolated quantities, but rather the simultaneous comparison of complementary state descriptors.

In particular:

\begin{itemize}
    \item purity diagnoses mixedness and coherence loss,
    \item fidelity measures proximity to the reference Bell state,
    \item concurrence and tangle quantify bipartite entanglement,
    \item entanglement of formation quantifies entanglement resources,
    \item CHSH quantities diagnose observable nonlocality.
\end{itemize}

\medskip

This separation is essential for the interpretation of fiber-induced effects.

A reduction in fidelity may arise solely from coherent phase rotations without affecting entanglement. By contrast, reductions in concurrence, tangle, and entanglement of formation indicate genuine degradation of quantum correlations.

Similarly, nonzero concurrence does not necessarily imply CHSH violation, especially in sufficiently mixed states.

\medskip

The coexistence of these complementary quantities therefore provides a multidimensional characterization of the propagation regime, allowing coherent rotations, statistical decoherence, entanglement degradation, and nonlocality suppression to be distinguished quantitatively throughout the simulator analysis.
